% =============================================================================
%  CHAPTER 3 — SYSTEM DESIGN
%  Generated 2026-09-23 from the implementation; every value below is read from
%  the code (see docs/methodology_facts.md).  Markers:
%     TODO   = prose you must write (kept as comments on purpose)
%     VERIFY = confirm the value before print
%     DECIDE = open decision that changes the text
%  Required packages: booktabs.  Optional: listings (verbatim prompt in App. A).
%
%  LABEL MAPPING: this chapter uses the sec:sd_* prefix (System Design);
%  sec:meth_* is reserved for the following Methodology chapter.
%  Structure is FLAT: no \subsection anywhere.  Run-in headings use \paragraph,
%  which keeps the table of contents flat while still giving the prose
%  landmarks (and they do not appear in the ToC at all).
%  chap:methodology now belongs to the following chapter; this one is
%  chap:system_design.
% =============================================================================

\chapter{System Design}
\label{chap:system_design}

% TODO (1 paragraph): state the two halves of the thesis artefact — the robot
% system (perception -> symbolic state -> decision -> Nav2) and the simulation
% suite — and say that this chapter documents both, while the measurement
% methodology is the subject of the next chapter.

\section{System Overview}
\label{sec:sd_overview}

\paragraph{Simulation platform.}

\begin{table}[htbp]
  \centering
  \small
  \caption{Simulation and compute platform.}
  \label{tab:platform}
  \begin{tabular}{ll}
    \toprule
    Component & Specification \\
    \midrule
    Host CPU & Intel Core i7-11800H (8 cores / 16 threads), capped at 2.4\,GHz \\
    Host GPU & NVIDIA GeForce RTX 3050\,Ti, 4\,GB VRAM \\
    Operating system & Ubuntu 22.04 container (\texttt{humble\_thesis\_engine}) \\
    Middleware & ROS 2 Humble \\
    Simulator & Ignition Fortress (Gazebo 6), physics step 1\,ms \\
    Simulated clock & Gazebo \texttt{/clock} relayed at 50\,Hz \\
    DDS & FastDDS, UDPv4-only transport, file-descriptor limit 65536 \\
    ROS domain & 31 \\
    Thermal precondition & run starts only at $\leq$ 60\,\textdegree C package temperature \\
    \bottomrule
  \end{tabular}
\end{table}

% TODO: one paragraph on WHY the clock is relayed at 50 Hz and why DDS is
% UDP-only (both were necessary for a stable multi-robot bring-up).
% DECIDE: the sensor rates and the clock relay are the current configuration;
% older runs used 10/5 Hz and an unrelayed clock.  State the canonical set here
% and keep it consistent with the era policy agreed for Chapter 4.

\paragraph{Robot platform.}

\begin{table}[htbp]
  \centering
  \small
  \caption{Simulated robot configuration (Leo Rover).}
  \label{tab:robot}
  \begin{tabular}{ll}
    \toprule
    Attribute & Value \\
    \midrule
    Kinematics & differential drive (URDF from \texttt{leo\_description}) \\
    Footprint & circular, radius 0.30\,m \\
    Max linear velocity & 0.5\,m/s \\
    Max angular velocity & 1.0\,rad/s \\
    Onboard sensors & 2D LiDAR ($\approx$10\,m range), RGB-D camera, bumper contact sensor \\
    LiDAR scan rate & 5\,Hz \\
    Camera rate / resolution & 3\,Hz, 320$\times$240 \\
    Command interface & \texttt{/\{ns\}/cmd\_vel} via the Gazebo bridge \\
    \bottomrule
  \end{tabular}
\end{table}

\paragraph{Behavioural conditions.}

% TODO (1-2 sentences): introduce that the same robot stack is evaluated under
% three decision layers, and that this is the independent variable of the study.

\begin{table}[htbp]
  \centering
  \small
  \caption{The three behavioural conditions. Everything except the decision layer
  is held constant.}
  \label{tab:conditions}
  \begin{tabular}{p{2.4cm}p{2.6cm}p{2.6cm}p{2.6cm}p{2.0cm}}
    \toprule
    Condition & Behaviour tree & Social layer & Human perception & LLM \\
    \midrule
    Nav2 (baseline) & stock \texttt{navigate\_to\_pose} with recoveries &
    none (geometric) & none & no \\
    BT (reactive) & \texttt{social\_nav.xml} & reactive yield and arc maneuvers &
    YOLO detector + rover proximity & no \\
    BT+LLM & \texttt{ai\_nav.xml} & reactive layer + agent decisions &
    same as BT & \texttt{gpt-4o-mini} \\
    \bottomrule
  \end{tabular}
\end{table}

% TODO (2-3 sentences): state explicitly that the baseline has no human
% perception at all, i.e. it is the "humans as costmap obstacles" condition
% described in Chapter 1 — this is deliberate and must be said, not implied.

\paragraph{Per-robot pipeline.}

% TODO (3-4 sentences walking through the pipeline):
%   Gazebo sensors -> ros\_gz bridges -> (a) lidar\_republish -> scan\_clean -> SLAM + costmaps
%                                     -> (b) YOLO detector -> symbolic human state
%   symbolic state + proximity -> decision layer (tree and/or agent) -> Nav2 controller
%   controller output -> cmd\_vel\_social -> velocity\_adaptor -> cmd\_vel -> wheels

\section{Perception and Symbolic States}
\label{sec:sd_perception}

\paragraph{LiDAR path.}

% TODO (2 sentences): Gazebo ignores the URDF frame id for the scan, so
% lidar\_republish rewrites it to \{ns\}/lidar\_link and republishes on
% /\{ns\}/lidar/scan\_clean, which is the single scan source for SLAM and for
% both costmaps.

\paragraph{Human detection.}

\begin{table}[htbp]
  \centering
  \small
  \caption{YOLO-based human detection parameters.}
  \label{tab:yolo}
  \begin{tabular}{ll}
    \toprule
    Parameter & Value \\
    \midrule
    Model & Ultralytics YOLOv8n (COCO class 0: person) \\
    Inference size & 640\,px \\
    Confidence threshold & 0.4 \\
    Inference rate limit & one frame per 0.33\,s ($\approx$3\,Hz) \\
    Range source & depth image at the box centre (5-px cross, then box median) \\
    Valid range & 0.2--50\,m \\
    Bearing & pinhole model; horizontal FOV 1.51844\,rad ($\approx$87\textdegree) \\
    Representative person & largest-area detection (nearest human) \\
    Close threshold & $0 < d < 1.5$\,m \\
    Detection hold & 1.0\,s on dropout \\
    Robot-as-person suppression & nearest rover $<$ 6.0\,m \emph{and} bearing difference $<$ 40\textdegree \\
    \bottomrule
  \end{tabular}
\end{table}

% TODO (1 paragraph): justify the suppression heuristic — a rover's body is
% detected as a person, and without suppression the social layer reacts to its
% own fleet as if it were pedestrians.

\paragraph{Symbolic state extraction.}

\begin{table}[htbp]
  \centering
  \small
  \caption{Symbolic state published per robot and consumed by the decision layer.}
  \label{tab:symbolic}
  \begin{tabular}{lll}
    \toprule
    State & Definition & Source \\
    \midrule
    \texttt{human\_detected} & a person is currently in view & YOLO detector \\
    \texttt{human\_close} & person within 1.5\,m & YOLO detector \\
    \texttt{human\_distance}, \texttt{human\_angle} & range and signed bearing to the nearest person (+ = left) & YOLO detector \\
    \texttt{robot\_close} & another rover within the interaction cone (30\textdegree) & proximity layer \\
    \texttt{robot\_angle}, \texttt{nearest\_robot\_dist}, \texttt{robot\_dca} & bearing, range and distance of closest approach to the nearest rover & proximity layer \\
    \texttt{obstacle\_min} & nearest LiDAR return & scan \\
    \texttt{speed\_scale} & social multiplier applied to the controller output & velocity adaptor \\
    \bottomrule
  \end{tabular}
\end{table}

% TODO (1 paragraph): state how these states are refreshed and logged, and that
% the decision layer never sees raw sensor data — only this symbolic state.

\section{Behavior Tree Design}
\label{sec:sd_bt_design}

\begin{verbatim}
RecoveryNode(retries=6) NavigateRecovery
+- PipelineSequence NavigateWithReplanning        (correct 1 Hz replanning)
   +- RateController(1.0 Hz)
   |  +- RecoveryNode(retries=1) -> ComputePathToPose | ClearEntireCostmap
   +- RecoveryNode(retries=1) SocialFollow
      +- ReactiveFallback HumanAwareFollow
         +- Timeout(4000 ms)  -> yield to a human, give up after 4 s
         +- ReactiveFallback RobotAwareFollow
            +- Timeout(10000 ms) -> yield to a rover, give up after 10 s
            +- ReactiveFallback ArcOrFollow
               +- Arc(adaptive, 0.3 m/s, 6 s, direction=away) -> FollowPath
\end{verbatim}

% TODO: a real diagram (TikZ) is a nicer figure than this listing — worth one
% afternoon if time allows.

\begin{table}[htbp]
  \centering
  \small
  \caption{Timing and retry parameters of the social tree.}
  \label{tab:bt_timing}
  \begin{tabular}{ll}
    \toprule
    Parameter & Value \\
    \midrule
    Navigation recovery retries & 6 \\
    Global replan rate & 1\,Hz \\
    Human yield bound & 4\,s \\
    Rover yield bound & 10\,s \\
    Escape maneuver & constant-curvature arc, 0.3\,m/s, 6\,s, away from the blocker \\
    In-place rotations & none (adhesive wheels; the robot always keeps rolling) \\
    \bottomrule
  \end{tabular}
\end{table}

% TODO (2 paragraphs): (a) why a PipelineSequence plus RateController rather
% than a SequenceStar root — a SequenceStar root re-ticked ComputePathToPose on
% every tick and caused planner timeouts; (b) why each ReactiveFallback holds at
% most two asynchronous children — with three, the yield branches stopped being
% re-ticked while FollowPath was active.

\section{Integration with Nav2}
\label{sec:sd_nav2_integration}

% TODO (2 paragraphs): the social layer never talks to the planner directly.
% It regulates motion through (i) the behaviour tree, and (ii) the velocity
% chain: the controller publishes cmd\_vel\_social, the velocity adaptor scales
% it by the social terms, and the result is published as cmd\_vel for the bridge.
% INCLUDE: progress checker relaxed to 0.25 m per 30 s because a deliberately
% pausing social robot would otherwise be declared non-progressing.

\section{LLM Social Layer}
\label{sec:sd_llm}

\begin{table}[htbp]
  \centering
  \small
  \caption{Agent configuration.}
  \label{tab:llm}
  \begin{tabular}{ll}
    \toprule
    Item & Value \\
    \midrule
    Model & \texttt{gpt-4o-mini} (hosted, accessed over HTTPS) \\
    Action vocabulary & \texttt{proceed}, \texttt{slow}, \texttt{yield}, \texttt{arc\_left}, \texttt{arc\_right}, \texttt{stop}, \texttt{back} \\
    Invocation policy & event-driven: one call per material change of the situation signature \\
    Input & symbolic situation report (nearest human and rover, free space, right-of-way) \\
    Output contract & single JSON object: \texttt{\{"action": ..., "reason": ...\}} \\
    Role & consultant — the tree executes the proposed action and the reactive layer remains available \\
    Logging & per-call action, reason, token counts and cost; per-robot totals \\
    Measured cost & \$0.011--0.013 per rover per 300 s run \\
    \bottomrule
  \end{tabular}
\end{table}

% TODO (3-4 bullets, the encoded social rules): humans have priority; the
% right-of-way is locked for the whole encounter (right-hand traffic); a rover
% beside or behind is already passed and must not be yielded to; adaptive
% passing arcs only when there is room; straight reverse is reserved for a
% genuine deadlock with a clear rear; escape from a stuck pass is allowed only
% for the yielding robot.
% NOTE: the full prompt is reproduced verbatim in Appendix A.

\section{Multi-Robot Extension}
\label{sec:sd_multi_robot}

% TODO (3 paragraphs):
%  (a) per-robot stacks: each rover runs its own SLAM and Nav2 stack, publishes
%      its own map and map->odom transform, and plans in its own map frame; maps
%      are never merged, because merging would couple every robot to a single
%      point of failure and is irrelevant to a purely local social layer.
%  (b) coordination is local only: rovers exchange no plans or goals, only what
%      each can perceive of the other (proximity layer) — the LLM right-of-way
%      rule is what makes the interaction predictable.
%  (c) scaling factor: the fleet size is varied (see Chapter 4) while each robot
%      keeps an independent stack.

\begin{table}[htbp]
  \centering
  \small
  \caption{Arena and scenario family.}
  \label{tab:scenarios}
  \begin{tabular}{ll}
    \toprule
    Item & Value \\
    \midrule
    Arena & 14 $\times$ 14\,m, walled on all sides \\
    Coverage grid & 14 $\times$ 14 cells of 1\,m (196 cells) \\
    Rover spawns & distinct corners, each facing the arena centre \\
    Humans & scripted walking actors on fixed lanes; lane span $\leq$ 7\,m; speed $\leq$ 0.5\,m/s \\
    Human lanes & two vertical lanes at $x = \pm 3.5$\,m, one shuttle lane at $y = 3$\,m \\
    Scenario matrix & \{2, 3, 4, 5\} rovers $\times$ \{1, 2, 3\} humans \\
    Obstacle variants & boxes and cylinders added to the 3- and 5-rover layouts \\
    \bottomrule
  \end{tabular}
\end{table}

% VERIFY: the results tree also contains 1-rover scenarios (needed as the
% single-robot end of RQ3); confirm how those were generated and add them to the
% matrix here if they are part of the study.
% TODO (contribution wording): the suite currently contains walking humans and
% static obstacles — it does NOT contain static humans. Either implement a
% standing pedestrian or reword contribution 4.

\section{Design Decisions and Alternatives Considered}
\label{sec:sd_rationale}

% TODO: one short paragraph per decision, each naming the rejected alternative:
%  - behaviour tree over a finite-state machine (modularity, readability,
%    reusability; Nav2 executes it natively)
%  - per-robot maps over a shared map (isolation, no single point of failure)
%  - event-driven LLM invocation over per-tick invocation (cost, latency, and
%    the fact that decisions are needed only at situation changes)
%  - agent as consultant over agent as controller (safety: the reactive tree
%    keeps working when the model is unavailable or wrong)
%  - nearest-person scene representation over full multi-person tracking
%    (sufficient for a local, reactive social layer; tracking is future work)
%  - bounded yields (4 s / 10 s) over indefinite yielding (the freezing robot
%    problem is precisely an unbounded yield)
%  - continuous arc escape over in-place rotation (adhesion, passenger comfort)

% =============================================================================
%  APPENDIX A — verbatim decision prompt (required for reproducibility)
% =============================================================================
\section*{Appendix A: Decision-layer prompt (verbatim)}
\label{app:prompt}

% The prompt below is the exact text sent to the model; it is part of the
% method, not an implementation detail, and must be reproduced unedited.

\begin{verbatim}
You are the social decision layer of a mobile robot navigating around other
robots and humans. Choose exactly ONE action:
- "proceed": continue at full speed along the current path (default when there
  is no conflict).
- "slow": keep driving but at reduced speed ...
- "yield": stop and let the other pass. Yield to a robot ONLY when the situation
  report says YOU are the YIELDING robot in the encounter; also yield to a human
  that is CLOSE (within 1.5 m) and ahead of you or crossing your path. NEVER
  yield when you HAVE the right-of-way. ...
- "arc_left" / "arc_right": a smooth curved detour to that side, then continue
  toward the goal (NEVER an in-place spin). ...
- "stop": emergency stop (e.g., collision imminent).
- "back": reverse STRAIGHT backwards (no turning) ONLY to break a DEADLOCK ...
Rules: humans always have priority ...
ADAPTIVE PASSING (robot-robot): ...
RIGHT-OF-WAY IS LOCKED FOR THE WHOLE ENCOUNTER (right-hand traffic ...)
STALL ESCAPE: ONLY the YIELDING robot may break a stuck pass. ...
Respond with a single JSON object, no markdown, no extra text:
{"action": "proceed|yield|slow|arc_left|arc_right|stop", "reason": "short reason"}
\end{verbatim}

% NOTE: the block above is abridged for reading. Copy the full text verbatim
% from swarm_basics/swarm_basics/llm_planner.py (SYSTEM_PROMPT) so the appendix
% is exact — an examiner may diff it against the implementation.
